\DocumentMetadata{
  pdfversion=2.0,
  pdfstandard=ua-2,
  lang=en-US,
  tagging=on,
  tagging-setup={math/setup=mathml-SE,tabsorder=structure}
}
\documentclass[11pt,letterpaper]{article}
\usepackage[margin=0.68in,includefoot]{geometry}
\usepackage{newcomputermodern}
\usepackage{amsmath,amscd,mathtools}
\providecommand{\square}{\mdlgwhtsquare}
\usepackage[hidelinks]{hyperref}
\hypersetup{
  pdftitle={Algebra PhD Exam, August 2024},
  pdfauthor={University of Florida Department of Mathematics},
  pdfsubject={Graduate mathematics examination},
  pdfdisplaydoctitle=true,
  bookmarksopen=true
}
\setcounter{secnumdepth}{0}
\setlength{\parindent}{0pt}
\setlength{\parskip}{0.28em}
\setlength{\emergencystretch}{3em}
\setlength{\leftmargini}{2.15em}
\setlength{\leftmarginii}{2.65em}
\setlength{\leftmarginiii}{3em}
\renewcommand{\labelenumi}{\textbf{\theenumi.}}
\pagestyle{empty}
\raggedbottom
\allowdisplaybreaks
\newenvironment{examproblems}[1]{%
  \begin{enumerate}%
  \setcounter{enumi}{#1}%
  \setlength{\topsep}{0.35em}%
  \setlength{\partopsep}{0pt}%
  \setlength{\itemsep}{0.48em}%
  \setlength{\parsep}{0.18em}%
}{\end{enumerate}}
\newenvironment{examsubparts}{%
  \begin{enumerate}%
  \setlength{\topsep}{0.18em}%
  \setlength{\partopsep}{0pt}%
  \setlength{\itemsep}{0.16em}%
  \setlength{\parsep}{0.1em}%
}{\end{enumerate}}
\newenvironment{examinstructions}{%
  \par\begingroup\itshape
}{%
  \par\endgroup\vspace{0.18em}
}
\makeatletter
\renewcommand\section{\@startsection{section}{1}{0pt}%
  {-1.25ex plus -.25ex minus -.1ex}%
  {0.55ex plus .1ex}%
  {\normalfont\large\bfseries}}
\renewcommand{\@maketitle}{%
  \newpage\null\vskip -1.2em%
  {\footnotesize\bfseries
    \href{https://www.ufl.edu/}{University of Florida}, \href{https://math.ufl.edu/}{Department of Mathematics}\par}%
  \vskip 0.18em%
  \tagstructbegin{tag=Title}%
    {\LARGE\bfseries\href{https://gma.math.ufl.edu/exams/algebra-phd/}{Algebra PhD Exam}, August 2024\par}%
  \tagstructend%
  \vskip 0.35em%
  \tagmcbegin{artifact}%
    \hrule height 1pt%
  \tagmcend%
  \vskip 0.55em%
}
\makeatother
\title{Algebra PhD Exam, August 2024}
\author{}
\date{}
\begin{document}
\maketitle
\thispagestyle{empty}
\begin{examinstructions}
Answer seven problems. If more than seven problems are answered only the first seven will be graded. Write your answers clearly in complete English sentences.

\par
Results from lectures or textbooks may be used without proof (within reason – don’t state results equivalent to the problem), but must be clearly stated.
\end{examinstructions}
\begin{examproblems}{0}
\item
Suppose \(K/k\) is a finite normal extension and let~\(G\) be the group of automorphisms of~\(K\) fixing every element of~\(k\). Denote by \(K^G\) the set of elements of~\(K\) fixed by every element of~\(G\). Show that \(K^G/k\) is a purely inseparable extension.
\item
Let~\(k\) be a field and \(A,B\) finitely generated~\(k\)-algebras. Let \(f:A\to B\) be a~\(k\)-algebra homomorphism. Show that if \(\mathfrak m\subset B\) is a maximal ideal, \(f^{-1}(\mathfrak m)\subset A\) is also a maximal ideal (Hint: Nullstellensatz). Recall that a~\(k\)-algebra~\(A\) is a ring homomorphism \(k\to A\) sending the identity to the identity, and if~\(A\) and~\(B\) are~\(k\)-algebras, a~\(k\)-algebra homomorphism \(A\to B\) is a commutative diagram 
\[
\begin{array}{ccc}
A & \longrightarrow & B \\
\uparrow & \nearrow & \\
k & &
\end{array}
\]
\item
This problem has three parts.
\begin{examsubparts}
\item[\textnormal{(i)}]
Let \(\mathcal C\) be a category. Define what it means for an object~\(X\) of \(\mathcal C\) to be a product of a set of objects \(\{X_i\}_{i\in I}\).
\item[\textnormal{(ii)}]
Let \(\mathcal T\) be the category of torsion abelian groups. Show that if \(\{X_i\}_{i\in I}\) is any set of objects of \(\mathcal T\), the product as defined in (i) exists.
\item[\textnormal{(iii)}]
Give an example of a set \(\{X_i\}_{i\in I}\) of objects of \(\mathcal T\) whose product is not isomorphic to the usual product (i.e. in the category of abelian groups).
\end{examsubparts}
\item
This problem has three parts.
\begin{examsubparts}
\item[\textnormal{(i)}]
Define what it means for a finite group~\(G\) to be nilpotent.
\item[\textnormal{(ii)}]
Show that if~\(p\) is a prime and~\(G\) is a~\(p\)-group then~\(G\) is nilpotent.
\item[\textnormal{(iii)}]
Show that \(A_4\) is solvable but not nilpotent.
\end{examsubparts}
\item
Suppose~\(R\) is an integral domain with fraction field~\(K\) and \(I,J\) are invertible fractional ideals of~\(R\). Recall that the product \(IJ\) is the~\(R\)-submodule of~\(K\) generated by the products \(ab\) with \(a\in I\) and \(b\in J\). Show that \(I\otimes_R J\simeq IJ\) as~\(R\)-modules.
\item
Suppose~\(R\) is a commutative ring, \(S\subseteq R\) is a multiplicative system (i.e. a nonempty subset closed under multiplication) and 
\[
0\to L\to M\to N\to 0
\]
 is an exact sequence of~\(R\)-modules. Show that the sequence 
\[
0\to S^{-1}L\to S^{-1}M\to S^{-1}N\to 0
\]
 of~\(R\)-modules is exact.
\item
For any set~\(S\) let \(F(S)\) denote the free group on~\(S\). Which sets~\(S\) have the property that \(F(S)\) is an abelian group? Give a proof of your answer from first principles.
\item
Suppose~\(R\) is an integral domain. Show that the polynomial ring \(R[X]\) is a Dedekind domain if and only if~\(R\) is a field.
\item
Suppose~\(R\) is a noetherian ring and~\(M\) is a finitely generated~\(R\)-module. Show that~\(M\) is noetherian.
\item
Let~\(R\) be a ring with identity. Show that any projective~\(R\)-module is flat.
\item
Let~\(p\) be a prime. Classify up to isomorphism all semisimple rings of cardinality \(p^6\).
\end{examproblems}
\end{document}
